\documentclass[11pt]{amsart}
\usepackage{geometry}                % See geometry.pdf to learn the layout options. There are lots.
\geometry{letterpaper}                   % ... or a4paper or a5paper or ... 
%\geometry{landscape}                % Activate for for rotated page geometry
%\usepackage[parfill]{parskip}    % Activate to begin paragraphs with an empty line rather than an indent
\usepackage{graphicx}
\usepackage{amssymb}
\usepackage{epstopdf}
\DeclareGraphicsRule{.tif}{png}{.png}{`convert #1 `dirname #1`/`basename #1 .tif`.png}

\title{Problema 4.4}
%\date{}                                           % Activate to display a given date or no date

\begin{document}
\maketitle
%\section{}
%\subsection{}
\noindent  Construya un diagrama de flujo para almacenar en un arreglo unidimensional los primeros 30 números primos. Al final imprima el arreglo correspondiente. \newline

\noindent  \textbf{Explicación de las variables} \newline

\noindent PRIMO: \hspace{0.5cm} Arreglo unidimensional de tipo entero. Se utiliza para almacenar los primeros 30 números primos.\newline
\noindent K: \hspace{0.5cm} Variable de tipo entero. Se usa como índice del arreglo y como va­riable de control de un ciclo.\newline
\noindent NUM: \hspace{0.5cm} Variable de tipo entero. Almacena el número sobre el que hay que determinar si es o no un número primo.\newline
\noindent DIVI: \hspace{0.5cm} Variable de tipo entero. Representa los posibles divisores del número que se esta analizando\newline
\end{document}  